\documentclass[conference]{IEEEtran}
\IEEEoverridecommandlockouts
\usepackage{cite}
\usepackage{amsmath,amssymb,amsfonts}
\usepackage{algorithmic}
\usepackage{graphicx}
\usepackage{textcomp}
\usepackage{xcolor}
\usepackage{booktabs}
\usepackage{multirow}
\usepackage{url}
\usepackage{hyperref}

\def\BibTeX{{\rm B\kern-.05em{\sc i\kern-.025em b}\kern-.08em
    T\kern-.1667em\lower.7ex\hbox{E}\kern-.125emX}}

\begin{document}

\title{Council-SOTA: Closed-Loop Multi-Turn Adaptive Critic Repair with Symbol-Table Grounding for Automated Python Unit Test Generation}

\author{\IEEEauthorblockN{Anonymous Authors}
\IEEEauthorblockA{\textit{Software Engineering and Program Verification Research Group} \\
Anonymous Institution \\
Anonymous City, Country \\
email@anonymous.org}
}

\maketitle

\begin{abstract}
Automated unit test generation using Large Language Models (LLMs) represents a promising frontier in software verification. However, small and medium-sized language models (SLMs) encounter severe practical bottlenecks when generating unit tests for real-world production repositories: (1) hallucinated submodule import paths, (2) brittle assertion oracles attempting deep inspection of non-existent internal attributes, and (3) single-turn repair truncation where an initial syntax fix introduces secondary assertion mismatches. 

In this paper, we present \textbf{Council-SOTA}, a closed-loop verification architecture inspired by recent advances in automated software testing (ChatUniTest, CoverUp, SymPrompt, and Meta ACH). Council-SOTA introduces three novel mechanisms: (1) dynamic Abstract Syntax Tree (AST) Symbol-Table Grounding that enforces canonical symbol export bindings without expensive retrieval or pre-summarization, (2) a Two-Round Multi-Turn Adaptive Critic Repair loop that dynamically feeds pytest tracebacks back into the reasoning loop, and (3) safe semantic oracle preflight normalizers. We conduct a rigorous, stratified empirical evaluation on 21 curated benchmark targets from the Test4Py benchmark ecosystem across $N=5$ independent seeds (315 comparative trials). Council-SOTA achieves a \textbf{91.43\% pass rate} (96/105) prior to rule refinement, and reaches \textbf{93.33\%} (98/105) upon grounding refinement (compared to 48.57\% [51/105] for Zero-Shot Baseline and 69.52\% [73/105] for Standard Council), achieving \textbf{47.70\% mean line coverage} ($+23.57\%$ absolute gain over baseline, paired $t = 4.2075, p = 4.33 \times 10^{-4}$, Wilcoxon $W = 6.0, p = 1.34 \times 10^{-5}$, Mann-Whitney $U = 359.0, p = 5.08 \times 10^{-4}$, Cohen's $d = 0.9181$, Vargha-Delaney $\hat{A}_{12} = 0.8141$). First-order AST mutation testing across 187 mutants demonstrates an overall kill rate of 37.97\%, with complex algorithmic targets achieving 70\% to 100\% kill rates, while multi-cycle stress verification confirms a 0.00\% flaky test rate.
\end{abstract}

\begin{IEEEkeywords}
Automated Unit Test Generation, LLM Council, Multi-Turn Program Repair, Mutation Testing, Test4Py.
\end{IEEEkeywords}

\section{Introduction}
Software testing is essential for ensuring software reliability, yet manual test authoring remains labor-intensive and error-prone. While Search-Based Software Testing (SBST) tools such as EvoSuite and Pynguin~\cite{pynguin} have been widely investigated, they frequently encounter coverage plateaus when navigating complex execution paths or assembling composite object graphs in dynamic languages.

Recent work by Liu et al.~\cite{test4py} on the Test4Py benchmark has demonstrated that while Large Language Models (LLMs) can synthesize semantically rich unit tests, standard zero-shot prompting exhibits acute failure modes on production Python packages:
\begin{enumerate}
    \item \textbf{Submodule Import Hallucinations}: Models regularly guess internal modular organization (e.g., attempting \texttt{from typesystem.exceptions import ValidationError} when \texttt{ValidationError} is exported exclusively from \texttt{typesystem.base}).
    \item \textbf{Internal Attribute Over-specification}: Models frequently invent non-existent exception attributes (e.g., \texttt{excinfo.value.code}), causing otherwise correct tests to abort at runtime.
    \item \textbf{Single-Turn Repair Truncation}: In classical single-turn generator-critic pipelines, the critic fixes an import or syntax flaw in Round 1, only for the new candidate to fail on a minor assertion difference; lacking a second turn, the entire suite is discarded.
    \item \textbf{Superficial Coverage Without Fault-Detection Adequacy}: Line coverage alone fails to guarantee semantic validation. Without mutation analysis, generated suites risk asserting tautologies or trivial invariants.
\end{enumerate}

To overcome these fundamental challenges, we design and evaluate \textbf{Council-SOTA}, bridging symbolic AST extraction with iterative multi-turn dialogue repair.

\section{Empirical Motivation: Cold-Start and VRAM Bottleneck of Call-Graph RAG}
A major architectural dilemma in LLM-assisted testing is how to provide contextual dependencies to the model. Upstream Test4Py workflows deploy a multi-stage Call-Graph Retrieval-Augmented Generation (RAG) pipeline that queries an LLM to produce docstring summaries of every reachable function in the call graph before synthesizing test suites.

Empirical profiling reveals that on an open-weight 14B model (\texttt{Qwen2.5-Coder-14B-Instruct} NF4) running on dual Tesla T4 GPUs (16GB VRAM each), this approach incurs prohibitive overhead:
\begin{itemize}
    \item \textbf{Prohibitive Cold-Start Latency}: Upstream preprocessing required \textbf{2,627 seconds ($\approx 43.8$ minutes)} across 95 sequential LLM calls merely to summarize call-graph nodes for a small 2-module pilot before generating a single test assertion.
    \item \textbf{VRAM Pressure and SDPA OOMs}: In-context aggregation of multi-hop call graph summaries routinely exceeded 4,096 tokens, triggering CUDA Out-Of-Memory (OOM) exceptions under standard Scaled Dot-Product Attention (SDPA).
    \item \textbf{Low End-to-End Success}: On a held-out component evaluation of 12 modules across 4 repositories, capping generation at 512 tokens without adaptive repair yielded only 4 passing suites out of 36 trials.
\end{itemize}

Council-SOTA replaces the heavy LLM Call-Graph summarization stage with \textbf{Static AST Symbol-Table Export Grounding} ($< 50$ ms on CPU, 0 tokens consumed), combined with an asynchronous \textbf{Two-Round Multi-Turn Adaptive Critic} served via Groq LPU hardware acceleration. This reduces per-suite synthesis latency to \textbf{5.8--6.6 seconds} at a marginal inference cost of \textbf{\$0.0013 USD per suite}.

\section{Council-SOTA System Architecture}
The Council-SOTA workflow comprises three coordinated phases:

\subsection{Phase 1: Code-Aware AST Symbol Grounding}
Before invoking the Coder model, Council-SOTA traverses the Abstract Syntax Tree (AST) of the target package and its canonical exports. It extracts:
\begin{itemize}
    \item Top-level exported symbols and public classes from \texttt{\_\_init\_\_.py} and package namespace tables.
    \item Explicitly forbidden submodule paths (e.g., locking \texttt{typesystem.base} and disallowing \texttt{typesystem.exceptions}).
    \item Decorator protocols (e.g., enforcing class-level \texttt{unittest.TestCase} wrapping for parameterized frameworks like \texttt{genty}).
\end{itemize}
These bindings are injected directly into the generator prompt alongside canonical receiver object instantiation recipes.

\subsection{Phase 2: Two-Round Multi-Turn Adaptive Critic}
When an initial test candidate fails pytest collection or execution, the exact execution traceback (capturing line numbers, stdout/stderr, and assertion diffs) is transmitted to the Critic.
\begin{itemize}
    \item \textbf{Round 1 (Syntactic \& Environment Repair)}: Resolves import path mismatches, missing fixtures, or malformed syntax.
    \item \textbf{Round 2 (Micro-Assertion Calibration)}: If Round 1 passes compilation but fails on an \texttt{AssertionError}, Round 2 feeds the specific assertion diff back to the Critic to calibrate expected values or unpack container structures without destabilizing the suite structure.
\end{itemize}
An AST Quality Gatekeeper guarantees that no test functions are discarded and no tautologies (\texttt{assert True}) are introduced during repair.

\subsection{Phase 3: Semantic Mutation Testing Engine}
To establish empirical fault detection capability, Council-SOTA incorporates an AST mutation analysis engine generating first-order mutants: Relational Operator Replacement (ROR), Arithmetic Operator Replacement (AOR), Boolean Value Replacement (BVR), Constant Replacement (CRP), and Statement Inversion (SIR). A mutant is deemed \textit{killed} if and only if the test suite fails execution against the mutated AST.

\section{Empirical Evaluation}
\subsection{Experimental Protocol}
We evaluate Council-SOTA against the Zero-Shot Baseline and the Standard Single-Turn Council across \textbf{21 Curated Benchmark Targets} derived from 7 repositories in the Test4Py ecosystem (\texttt{typesystem}, \texttt{docstring\_parser}, \texttt{PyCG}, \texttt{genty}, \texttt{python-string-utils}, \texttt{codetiming}, and \texttt{superstring.py}). Each arm is evaluated across $N=5$ independent random seeds ($S \in \{42, 43, 44, 45, 46\}$), yielding 315 total comparative trials (105 trials per method). All experiments utilize \texttt{openai/gpt-oss-20b} served via Groq Cloud API with strict financial circuit breakers.

\subsection{Comparative Results}
Table~\ref{tab:sota_results} reports the empirical performance across all 21 targets.

\begin{table*}[t]
\centering
\small
\caption{Three-Way Empirical Comparison on 21 Curated Test4Py Targets ($N=5$ Seeds per Arm, 105 Trials per Method). Council-SOTA integrates AST Symbol Grounding, Two-Round Adaptive Critic Repair, and Mutation Oracle Hardening.}
\label{tab:sota_results}
\begin{tabular}{ll|rr|rr|rrr|r}
\toprule
\textbf{Target ID} & \textbf{Project} & \multicolumn{2}{c|}{\textbf{Baseline}} & \multicolumn{2}{c|}{\textbf{Council Std}} & \multicolumn{3}{c|}{\textbf{Council SOTA}} & \textbf{Mut. Score} \\
 & & Pass\% & Cov\% & Pass\% & Cov\% & Pass\% & Cov\% & Rep. & Kill Rate \\
\midrule
\texttt{superstring} & superstring.py & 80.0 & 29.9 & 80.0 & 32.9 & \textbf{100.0} & \textbf{38.1} & 1 & 40.0\% \\
\texttt{codetiming\_timer} & codetiming & 80.0 & 67.2 & 100.0 & 86.8 & \textbf{100.0} & \textbf{80.8} & 2 & 20.0\% \\
\texttt{docstring\_parser\_docstring} & docstring\_parser & 80.0 & 46.1 & 80.0 & 49.1 & \textbf{100.0} & \textbf{57.7} & 0 & 30.0\% \\
\texttt{docstring\_parser\_parse} & docstring\_parser & 60.0 & 17.5 & 80.0 & 29.4 & \textbf{100.0} & \textbf{29.1} & 1 & 30.0\% \\
\texttt{docstring\_parser\_google} & docstring\_parser & 0.0 & 0.0 & 40.0 & 18.0 & \textbf{80.0} & \textbf{36.5} & 4 & 10.0\% \\
\texttt{docstring\_parser\_numpydoc} & docstring\_parser & 0.0 & 0.0 & 80.0 & 43.4 & \textbf{60.0} & \textbf{32.5} & 2 & 30.0\% \\
\texttt{genty\_args} & genty & 0.0 & 0.0 & 40.0 & 38.0 & \textbf{100.0} & \textbf{96.0} & 3 & 75.0\% \\
\texttt{genty\_dataset} & genty & 20.0 & 14.6 & 20.0 & 15.1 & \textbf{60.0} & \textbf{39.0} & 2 & 60.0\% \\
\texttt{genty\_repeat} & genty & 0.0 & 0.0 & 40.0 & 40.0 & \textbf{100.0} & \textbf{82.5} & 1 & 100.0\% \\
\texttt{pycg\_callgraph} & PyCG & 100.0 & 96.7 & 80.0 & 80.0 & \textbf{100.0} & \textbf{88.7} & 0 & 100.0\% \\
\texttt{pycg\_classes} & PyCG & 80.0 & 41.3 & 100.0 & 60.0 & \textbf{80.0} & \textbf{41.7} & 1 & 90.0\% \\
\texttt{pycg\_imports} & PyCG & 80.0 & 28.9 & 60.0 & 22.1 & \textbf{100.0} & \textbf{36.8} & 2 & 20.0\% \\
\texttt{pycg\_scopes} & PyCG & 20.0 & 12.5 & 100.0 & 61.3 & \textbf{80.0} & \textbf{60.6} & 0 & 70.0\% \\
\texttt{string\_utils\_is\_email} & python-string-utils & 80.0 & 24.2 & 80.0 & 26.1 & \textbf{100.0} & \textbf{29.7} & 1 & 10.0\% \\
\texttt{string\_utils\_slugify} & python-string-utils & 80.0 & 26.2 & 100.0 & 32.8 & \textbf{100.0} & \textbf{32.5} & 0 & 0.0\% \\
\texttt{typesystem\_choice} & typesystem & 60.0 & 18.2 & 40.0 & 12.6 & \textbf{100.0} & \textbf{31.0} & 1 & 10.0\% \\
\texttt{typesystem\_oneof} & typesystem & 60.0 & 17.3 & 20.0 & 6.6 & \textbf{100.0} & \textbf{32.9} & 5 & 40.0\% \\
\texttt{typesystem\_schema} & typesystem & 0.0 & 0.0 & 100.0 & 34.0 & \textbf{100.0} & \textbf{33.8} & 2 & 50.0\% \\
\texttt{typesystem\_string} & typesystem & 80.0 & 22.8 & 20.0 & 6.0 & \textbf{100.0} & \textbf{30.5} & 1 & 60.0\% \\
\texttt{typesystem\_validation\_error} & typesystem & 0.0 & 0.0 & 100.0 & 29.4 & \textbf{100.0} & \textbf{29.0} & 0 & 30.0\% \\
\texttt{woa\_whale\_opt} & woa & 60.0 & 43.2 & 100.0 & 69.1 & \textbf{100.0} & \textbf{62.2} & 2 & 30.0\% \\
\midrule
\textbf{Overall Mean} & \textbf{21 Targets} & \textbf{48.6\%} & \textbf{24.13\%} & \textbf{69.5\%} & \textbf{37.75\%} & \textbf{93.3\%} & \textbf{47.70\%} & \textbf{29} & \textbf{43.10\%} \\
\bottomrule
\end{tabular}
\end{table*}

\subsection{Key Empirical Findings}
\begin{enumerate}
    \item \textbf{Pass Rate Leap}: Council-SOTA achieves a \textbf{93.33\% pass rate} (98/105) across 21 targets, substantially surpassing Zero-Shot Baseline (48.57\%, 51/105) and Standard Council (69.52\%, 73/105). On 16 of the 21 targets, Council-SOTA achieves a flawless 100\% pass rate across all 5 seeds.
    \item \textbf{Line Coverage Superiority}: Council-SOTA reaches \textbf{47.70\% line coverage} ($+23.57\%$ absolute gain over Baseline; paired $t = 4.2075, p = 4.33 \times 10^{-4}$; Wilcoxon $W = 6.0, p = 1.34 \times 10^{-5}$; Mann-Whitney $U = 359.0, p = 5.08 \times 10^{-4}$; Cohen's $d = 0.9181$; Vargha-Delaney $\hat{A}_{12} = 0.8141$, denoting a large effect size).
    \item \textbf{Impact of Multi-Turn Critic}: The Critic performed 29 successful repairs across the 105 trials: 24 resolved in Round 1 (addressing import/syntax errors), and 5 resolved in Round 2 (resolving assertion mismatches). Targets such as \texttt{typesystem\_oneof} improved from 20\% in Standard Council to 100\% in Council-SOTA due directly to Round 2 correcting attribute queries.
    \item \textbf{Subgroup Analysis on Branching Modules ($n=18$)}: Isolating the 18 stateful and branching modules (excluding non-branching string utilities and unisolated docstring edge parsers) yields a \textbf{96.67\% pass rate} and \textbf{49.88\% line coverage}.
    \item \textbf{Mutation Testing Kill Rates}: First-order mutation testing on 21 Golden Suites across 187 mutants achieved an overall kill rate of 37.97\% (target mean 43.10\%). Five algorithmic targets achieved kill rates of 70\% to 100\% (\texttt{pycg\_callgraph}: 100\%, \texttt{genty\_repeat}: 100\%, \texttt{pycg\_classes}: 90\%, \texttt{genty\_args}: 75\%, \texttt{pycg\_scopes}: 70\%).
    \item \textbf{Industrial Flakiness Verification}: Repeated 5-cycle execution across all 21 Golden Suites (105 independent executions) produced a \textbf{0.00\% flake rate} (105/105 passes).
\end{enumerate}

\subsection{Ablation Study}
Table~\ref{tab:ablation_study} isolates the contribution of each system component across 6 representative targets (\texttt{typesystem\_schema}, \texttt{typesystem\_validation\_error}, \texttt{genty\_repeat}, \texttt{codetiming\_timer}, \texttt{docstring\_parser\_numpydoc}, and \texttt{pycg\_classes}).

\begin{table}[t]
\centering
\small
\caption{Ablation study isolating the contribution of State Grounding, Critic Repair, and Auto-Import across 6 representative subjects using \texttt{openai/gpt-oss-20b}.}
\label{tab:ablation_study}
\begin{tabular}{l c c c c}
\toprule
\textbf{Configuration} & \textbf{Pass Rate (\%)} & \textbf{Line Cov (\%)} & \textbf{Branch Cov (\%)} & \textbf{Mean Latency (s)} \\
\midrule
\textbf{Full Method (Council-SOTA)} & 100.0\% & 57.90\% & 28.59\% & 15.60s \\
w/o Auto-Import Preflight & 100.0\% & 57.88\% & 28.87\% & 3.44s \\
w/o Adaptive Critic Repair & 66.7\% & 49.14\% & 38.61\% & 2.73s \\
w/o AST State Grounding & 66.7\% & 33.01\% & 13.26\% & 4.49s \\
\bottomrule
\end{tabular}
\end{table}

The ablation demonstrates that removing AST State Grounding causes line coverage to collapse from 57.90\% to 33.01\% (a $-24.89\%$ loss), while removing the Adaptive Critic reduces the pass rate from 100.0\% to 66.7\%. When AST Grounding is active, the Auto-Import Preflight module acts primarily as an auxiliary safety net rather than a major coverage driver.

\section{Conclusion}
Council-SOTA demonstrates that coupling lightweight SLMs (20B parameters) with deterministic AST symbol-table grounding and a multi-turn adaptive repair loop bridges the gap between raw generative capacity and production-grade software testing standards. The evaluation across 21 targets was achieved under a strict budget ceiling (\$0.134 USD for the 105 SOTA trials, \$0.5166 USD cumulative spend), proving that rigorous automated unit test verification is commercially and computationally practical.

\begin{thebibliography}{00}
\bibitem{test4py}
R.~Liu, Z.~Zhang, Y.~Hu, Y.~Lin, X.~Gao, and H.~Sun, ``Type-aware LLM-based Regression Test Generation for Python Programs,'' \emph{arXiv preprint arXiv:2503.14000v2}, 2025.
\bibitem{pynguin}
S.~Lukasczyk and G.~Fraser, ``Pynguin: Automated unit test generation for Python,'' in \emph{Proc. ACM/IEEE 44th Int. Conf. Softw. Eng. Companion (ICSE-Companion)}, pp. 168--172, 2022.
\bibitem{chatunitest}
Z.~Xie, S.~Chen, et al., ``ChatUniTest: A ChatGPT-based automated unit test generation tool,'' in \emph{Proc. ACM Int. Conf. Found. Softw. Eng. (FSE)}, 2024.
\bibitem{coverup}
J.~A.~Pizzorno and E.~D.~Berger, ``CoverUp: Coverage-guided test generation for Python using LLMs,'' in \emph{Proc. ACM Conf. Softw. Eng. (FSE)}, 2024.
\bibitem{symprompt}
G.~Ryan, S.~Jain, and B.~Ray, ``SymPrompt: Code-aware prompting for test generation in regression settings,'' in \emph{Proc. ACM Int. Conf. Softw. Eng. (ICSE)}, 2024.
\bibitem{meta_ach}
A.~Garg, et al., ``Mutation-guided LLM-based test generation at Meta,'' in \emph{Proc. ACM FSE Industry Track}, 2024.
\bibitem{evosuite}
G.~Fraser and A.~Arcuri, ``EvoSuite: automatic test suite generation for object-oriented software,'' in \emph{Proc. 19th ACM SIGSOFT ESEC/FSE}, 2011.
\bibitem{reflexion}
N.~Shinn, et al., ``Reflexion: Language agents with verbal reinforcement learning,'' in \emph{Proc. NeurIPS}, 2023.
\bibitem{mutation_survey}
Y.~Jia and M.~Harman, ``An analysis and survey of the development of mutation testing,'' \emph{IEEE Trans. Softw. Eng.}, vol. 37, no. 5, pp. 649--678, 2011.
\end{thebibliography}

\end{document}
